\chapter{Soluciones de los ejercicios}\label{app:answers}
Estas respuestas corresponden a los cuatro ejercicios de estudio de cada capítulo. Desarrollan el razonamiento necesario y remiten a las derivaciones del cuerpo cuando éstas forman parte de la respuesta. Los casos M del manual y los problemas de los cuadernos tienen su propia solución completa junto al enunciado o en la parte separada de soluciones.
\section*{Capítulo 1}
\begin{enumerate}
\item El motor transporta reactantes y expulsa masa con momento opuesto al incremento del vehículo. No necesita apoyarse en aire. En el modelo descrito, eliminar la descarga eliminaría el flujo de momento; el vacío no suministra una fuerza que lo reemplace.
\item $I=F\Delta t=0.12(250)=30\,\mathrm{N\,s}$. Es transferencia de momento. Para convertirlo en velocidad se requiere masa y su evolución; sólo con masa prácticamente constante y sin otras fuerzas se aproxima $\Delta v=I/m$.
\item De $P_j=\dot mv_e^2/2$ resulta $\dot m_2/\dot m_1=1/4$ y $F_2/F_1=1/2$. Si no hay término de presión y toda la masa tiene esa velocidad, $I_{sp,2}/I_{sp,1}=2$.
\item Incluir fuente y conversión aumenta masa seca y puede disminuir la relación de masas útil. La alternativa que ahorra propulsante no tiene por qué minimizar masa inicial ni duración. Deben compararse fronteras completas de misión.
\end{enumerate}
\section*{Capítulo 2}
\begin{enumerate}
\item $[c^*]=[pA]/[\dot m]=\mathrm{m/s}$; $[C_F]=1$; $[I]=\mathrm{N\,s}$; $[I_{sp}]=\mathrm s$. Las dimensiones proceden de definiciones y no de cómo se denomine cada magnitud.
\item $\overline I_{sp}=(2\cdot240+5\cdot300)/7=282.857\,\mathrm s$. Se suman impulsos $g_0m_iI_{sp,i}$ y se divide por el peso convencional total consumido; la duración de fases no se ha dado ni hace falta para ese promedio.
\item El sistema consume también la descarga auxiliar. Se debe sumar su momento axial y su masa en $F_{tot}/(g_0\dot m_{tot})$. Usar sólo masa de cámara principal aumenta artificialmente el cociente si se mantiene empuje total.
\item Si $c=v_e+F_p/\dot m$, entonces $\dot mc=\dot mv_e+F_p$. Añadir después $F_p$ produciría $\dot mv_e+2F_p$, que no coincide con el balance original.
\end{enumerate}
\section*{Capítulo 3}
\begin{enumerate}
\item Con kg/kmol, $R=8314.462618/24=346.435942\,\mathrm{J/(kg\,K)}$. Con kg/mol, $R=8.314462618/0.024$, el mismo valor. Mezclar 8.314 con 24 introduce un factor mil.
\item El factor $1+(\gamma-1)M^2/2$ tiende a uno. Por tanto $T_0/T$, $p_0/p$ y $\rho_0/\rho$ tienden a uno para el modelo perfecto con $\gamma>1$ fijo.
\item $s_2-s_1=c_p\ln(T_2/T_1)-R\ln(p_2/p_1)$. Una temperatura menor sólo fija el signo del primer término; el cambio de presión decide el segundo. Isentropía exige suma nula, no simplemente enfriamiento.
\item Sí: un choque normal adiabático conserva $h_0=c_pT_0$ y genera entropía. Con $T_0$ constante, $\Delta s=-R\ln(p_{02}/p_{01})>0$ implica $p_{02}<p_{01}$.
\end{enumerate}
\section*{Capítulo 4}
\begin{enumerate}
\item Independencia elimina covarianza: $u_I/I=\sqrt{0.012^2+0.008^2}=0.0144222$, es decir, $1.44222\,\%$. Es incertidumbre estándar linealizada, sin factor de cobertura.
\item Se necesita saber qué pérdida representa cada factor y si afecta a caudal, rapidez, dirección o presión. Si la corrección de $C_F$ ya contiene la reducción cinética introducida en $v_e$, aplicarla de nuevo duplicaría el efecto.
\item Entre otros: Reynolds, Mach, Prandtl y razones de temperatura de pared y remanso. Con igual expansión geométrica, cambiar tamaño altera Reynolds si no se compensan propiedades y velocidad; no se garantiza semejanza viscosa.
\item Introducir un diámetro doble debe cuadruplicar $A_t=\pi d_t^2/4$ y el caudal crítico a estados totales fijos. Si el cálculo sólo duplica área o caudal, se ha usado una dependencia lineal incorrecta. Dimensionalmente, un diámetro sin cuadrado no tiene unidad de área.
\end{enumerate}
\section*{Capítulo 5}
\begin{enumerate}
\item Continuidad da $d\rho/\rho+du/u+dA/A=0$. Momento da $dp=-\rho udu$. La isentropía permite sustituir $dp=a^2d\rho$, obteniendo $d\rho/\rho=-M^2du/u$ y finalmente $dA/A=(M^2-1)du/u$.
\item $du/u\simeq-0.01/(0.4^2-1)=0.0119048$: aumento de $1.19048\,\%$. Es una linealización; el error de una variación finita contiene términos de segundo orden y cambios de Mach.
\item Para $M=2$, $du/u=(dA/A)/3$. Acelerar requiere $dA>0$; una expansión de 1\,\% daría aproximadamente $0.333333\,\%$ de incremento de rapidez.
\item Un divergente subsónico desacelera y recupera presión. Se necesitan estados de remanso, geometría anterior y contrapresión para seleccionar régimen y rama. La forma local por sí sola no determina Mach.
\end{enumerate}
\section*{Capítulo 6}
\begin{enumerate}
\item Para $\gamma=1.4$ y expansión tres, las raíces son aproximadamente $M_{sub}=0.197449$ y $M_{sup}=2.637416$. Sustituir cada una en la relación área-Mach recupera tres; el intervalo de búsqueda debe conservarse para distinguirlas.
\item El primer ahogo CD usa $M_{sub}$: $p_b/p_0\simeq0.973182$. La convergente simple alcanza su umbral en $(2/2.4)^{3.5}=0.528282$. La diferencia es recuperación subsónica en el divergente.
\item La expresión crítica es lineal en $A_t$, por lo que duplicarlo duplica caudal si el resto y ahogo se mantienen. Duplicar sólo salida mantiene caudal en ese régimen y cambia expansión y prestaciones.
\item $A^*_2/A^*_1=p_{01}/p_{02}=1/0.8=1.25$. A caudal y temperatura total constantes, una pérdida de presión total aumenta el área crítica de referencia un 25\,\%.
\end{enumerate}
\section*{Capítulo 7}
\begin{enumerate}
\item Se parte de $C_F=v_e/c^*+(r-p_a/p_c)\varepsilon$. Sustituir las expresiones energéticas de $v_e$ y de descarga para $c^*$ produce \eqref{eq:CF}. Multiplicar por $c^*$ devuelve $v_e+(p_e-p_a)A_e/\dot m=F/\dot m$.
\item $\Delta F=(80000-20000)(0.020)=1200\,\mathrm N$. El signo es positivo porque baja la presión exterior; se mantiene estado interno y compatibilidad de régimen.
\item Para Mach grande, $A\sim M^{2/(\gamma-1)}$ y $p\sim M^{-2\gamma/(\gamma-1)}$. Su producto escala como $M^{-2}$ y tiende a cero.
\item La geometría y cámara fijan $p_e$ en la rama ideal, mientras $p_a$ baja durante ascenso. La igualdad sólo se mantiene en un punto o si se modifica el estado o geometría de forma compatible.
\end{enumerate}
\section*{Capítulo 8}
\begin{enumerate}
\item Para $M_1=2$, $M_2=0.577350$, $p_2/p_1=4.5$, $\rho_2/\rho_1=2.666667$, $T_2/T_1=1.6875$ y $p_{02}/p_{01}=0.720874$. La energía total cierra y $\Delta s/R=-\ln(0.720874)>0$.
\item Para expansión tres y gamma 1.4: primer bloqueo $0.973182$, choque en salida aproximadamente $0.375959$ y adaptación $0.047299$. Son umbrales de un modelo ideal y no criterios de separación viscosa.
\item Con $T_{02}=T_{01}$, $\Delta s=-R\ln(p_{02}/p_{01})$. Si también se conservara $p_0$, resultaría entropía nula, incompatible con un choque compresivo de intensidad finita del modelo.
\item Se propone un área de choque, se calcula rama supersónica anterior, se aplica salto, se actualiza $A^*_2=A_t p_{01}/p_{02}$ y se sigue rama subsónica hasta salida. Se ajusta la posición hasta que la presión de salida iguale contrapresión.
\end{enumerate}
\section*{Capítulo 9}
\begin{enumerate}
\item $\lambda=(1+\cos12^\circ)/2=0.989074$. Corrige el flujo axial de momento del reparto supuesto; el término de presión requiere su propia geometría y distribución.
\item $\eta_n=C_v^2=0.9216$. Tratar ambos como 0.96 sobreestima el rendimiento cinético en 0.0384 absolutos. Para recuperar velocidad desde rendimiento energético se usa raíz cuadrada.
\item Deben coincidir reactantes, entalpías de entrada, cámara, presión o geometría final, especies y fases admitidas y punto de congelación. Después se comparan sólo los mecanismos químicos que cambian entre modelos.
\item Registrar tiempo, $p_c$, caudal total y auxiliares, $A_t$ inicial e instantánea, $A_e$, ambiente, empuje y sus incertidumbres. El cálculo debe indicar qué columna de área usa y no inferir eficiencia química sólo de $c^*$.
\end{enumerate}
\section*{Capítulo 10}
\begin{enumerate}
\item Si $r_b$ está en m/s y presión en Pa, $[a]=\mathrm{m\,s^{-1}Pa^{-n}}$. Como $p_{Pa}=10^6p_{MPa}$, $a_{MPa}=10^{6n}a_{Pa}$. La velocidad física no cambia.
\item $r_b/r_{ref}=1.20^{0.35}\simeq1.06589$. Se usa un cociente adimensional y un exponente local; no es una tasa absoluta.
\item La masa y densidad fijan volumen, pero no superficie activa ni su evolución. También faltan ley de regresión, garganta, estado del gas y ambiente para determinar empuje.
\item $\dot m_g=\int_{A_b}\rho_pr_b(\mathbf x)\,dA$. Se reduce a producto si densidad y regresión son uniformes sobre la superficie o si se utiliza una media superficial correctamente ponderada y definida.
\end{enumerate}
\section*{Capítulo 11}
\begin{enumerate}
\item $V_{cons}=4\pi[R_0^3-(R_0-w)^3]/3$. Derivar da $A_b=4\pi(R_0-w)^2$, decreciente para $0<w<R_0$.
\item A mitad de profundidad, $r/r_0=(5+1)/2=3$ y esa es la razón de áreas. A mitad de volumen, $r/r_0=\sqrt{(25+1)/2}=\sqrt{13}=3.60555$.
\item Sea $a=(1-2n)/(1-n)$. Como $(1+x)^a=1+a\ln(1+x)+O(a^2)$, el cociente $[(1+x)^a-1]/a$ tiende a $\ln(1+x)$.
\item Erosión de garganta, cambio de propiedades o regresión y variación de presión ambiente pueden modificar empuje aunque área activa sea constante. El almacenamiento y transitorios aportan otros mecanismos.
\end{enumerate}
\section*{Capítulo 12}
\begin{enumerate}
\item $p_2/p_1=1.15^{1/(1-0.4)}\simeq1.26230$. El caudal tiene el mismo cociente con garganta y $c^*$ fijos. La regresión tiene razón $(p_2/p_1)^{0.4}$ y cierra generación.
\item $A_{b,2}/A_{b,1}=(p_2/p_1)^{1-n}=1.08^{0.7}\simeq1.05535$ para garganta fija. No se invierte el exponente de esta relación inversa.
\item Derivar $m_g=pV/(RT)$ da $\dot m_g=V\dot p/(RT)+p\dot V/(RT)-pV\dot T/(RT^2)$. Despejar $V\dot p$ produce \eqref{eq:solidDynamic}. A masa y volumen constantes queda $\dot p/p=\dot T/T$.
\item La integración debe tender a $P_{eq}=a^{1/(1-n)}$ para el estado estable del modelo. Se comprueba presión positiva, tolerancia numérica y residuo $aP^n-P$. El registro de comprobaciones incluye el caso $a=1.02,n=0.4$.
\end{enumerate}
\section*{Capítulo 13}
\begin{enumerate}
\item $\tau_{rel}=\tau_0/(1-n)$. Para 0.30 es $1.42857\tau_0$ y para 0.80 es $5\tau_0$. Mayor exponente produce recuperación más lenta en el modelo concentrado.
\item La perturbación lineal de equilibrio es $0.1/0.6=0.166667$; la exacta es $1.1^{1/0.6}-1\simeq0.17216$. El error debe compararse en presión total o perturbación indicando el denominador empleado.
\item El modelo medio no contiene fase espacial ni energía de modos acústicos. Puede amortiguar una presión uniforme y omitir un modo que recibe calor en fase; son espacios de estados distintos.
\item Aislamiento introduce resistencia $\delta/k$; acumulación sensible introduce $mc\dot T$; ablación incorpora flujo másico y entalpía efectiva, además de frontera móvil. Se deben cerrar junto a las otras rutas térmicas.
\end{enumerate}
\section*{Capítulo 14}
\begin{enumerate}
\item Con $\mathcal R=2.8$, $\dot m_f=\dot m/3.8$ y $\dot m_o=2.8\dot m/3.8$. Suman el total y su cociente recupera 2.8.
\item Para masas iguales, $\rho_{eq}=2/(1/900+1/1200)=1028.571\,\mathrm{kg/m^3}$. El volumen se suma como masa/densidad y por ello no se obtiene la media aritmética 1050.
\item Si $\mathcal R=a+bt$ y $\dot m=M$ constante, $m_f=\int M/(1+a+bt)dt=(M/b)\ln[(1+a+bt_f)/(1+a)]$ para $b\ne0$. $m_o=Mt_f-m_f$. Para $b=0$ se toma el límite constante.
\item $c_{sist}=c_m+\beta(c_a-c_m)$. Su derivada respecto a beta es negativa si $c_a<c_m$; dividir por $g_0$ conserva esa tendencia para impulso específico.
\end{enumerate}
\section*{Capítulo 15}
\begin{enumerate}
\item En el orden H$_2$O,H$_2$,O$_2$,OH,O,H, las filas H y O son $(2,2,0,1,0,1)$ y $(1,0,2,1,1,0)$. Toda reacción admisible tiene vector estequiométrico en el núcleo de esa matriz, conservando ambos elementos.
\item TP fija temperatura en lugar de resolver energía adiabática. Para temperatura de cámara ideal se requiere un problema HP con entalpías de reactantes y presión, resolviendo composición y temperatura simultáneamente.
\item $3.0414\,\mathrm{kJ/(kg\,K)}=3041.4\,\mathrm{J/(kg\,K)}$. Usar 3.0414 en una ecuación en julios reduce ese término energético por mil.
\item Una segunda ejecución comprueba reproducibilidad local. Una implementación independiente contrasta cálculo. Un balance elemental comprueba conservación. Datos experimentales bajo condiciones conocidas permiten evaluar capacidad predictiva. No son etiquetas intercambiables.
\end{enumerate}
\section*{Capítulo 16}
\begin{enumerate}
\item $\tau_r=\rho_cV_c/\dot m$, con $\rho_c=p_c/(RT_c)$ y $\dot m=p_cA_t/c^*$, da $\tau_r=c^*L^*/(RT_c)$. Las dimensiones son $(\mathrm{m/s})\mathrm m/(\mathrm{m^2/s^2})=\mathrm s$.
\item $t_v=d_0^2/K$ se multiplica por cuatro, y la masa inicial proporcional a $d_0^3$ se multiplica por ocho. El mismo diámetro medio no representa automáticamente una distribución amplia de gotas.
\item Con $\Delta p=p_s-p_c$, $\delta\dot m/\dot m=-\delta p_c/(2\Delta p)$ para suministro fijo y coeficientes constantes. La subida de cámara disminuye entrada en la respuesta estática.
\item La recirculación, trayectorias y mezcla pueden producir tiempos muy distintos para partículas individuales. Igualar masa almacenada/caudal sólo fija una escala media; no fija la distribución ni su cola de residencia corta.
\end{enumerate}
\section*{Capítulo 17}
\begin{enumerate}
\item Igualar los tres flujos y sumar saltos da $T_{aw}-T_l=q''(1/h_g+\delta/k+1/h_l)$. Después $T_{wg}=T_{aw}-q''/h_g$ y $T_{wl}=T_l+q''/h_l$.
\item $\Delta T=24000/(0.80\cdot2500)=12\,\mathrm K$. La potencia recuperada $\dot mc_p\Delta T$ es exactamente 24 kW. No determina máximo térmico de pared.
\item Transferir calor con temperatura finita genera $dS_{gen}=\delta Q(1/T_c-1/T_h)>0$. Recuperarlo en una corriente no elimina esa irreversibilidad ni las pérdidas de presión.
\item Con ebullición se usa entalpía de fase y calidad; con almacenamiento se añade capacidad térmica por derivada temporal; con ablación entran masa saliente, entalpías y frontera móvil. No se aplica inalterado el balance estacionario de tres resistencias.
\end{enumerate}
\section*{Capítulo 18}
\begin{enumerate}
\item Al duplicar volumen, $p/p_0=0.5$ en el isotérmico y $2^{-1.4}=0.378929$ en el adiabático reversible. En el segundo, $T/T_0=2^{-0.4}$.
\item En el reservorio rígido, $d(mc_vT)=c_pT\,dm$, con $dm<0$. Expandir da $m c_vdT=RTdm$, y por integración $T/T_0=(m/m_0)^{\gamma-1}$. No hay aumento de volumen.
\item La presión de estado es absoluta. Se suma la presión exterior de referencia a la lectura manométrica, respetando unidades; la presión de vapor también debe expresarse sobre una referencia compatible.
\item Escriba $a=K(A_t/c^*)^2>0$. La raíz positiva de $ap_c^2+p_c-p_{dep}=0$ es $p_c=[-1+\sqrt{1+4ap_{dep}}]/(2a)$. La otra es negativa. La ecuación implica $p_{dep}=p_c+ap_c^2>p_c$.
\end{enumerate}
\section*{Capítulo 19}
\begin{enumerate}
\item $P_p=3(1.2\times10^6)/(1000\cdot0.75)=4800\,\mathrm W$. El trabajo hidráulico es 1200 J/kg y el de eje 1600 J/kg.
\item Multiplicar $\mathcal T=\dot m(r_2c_{\theta2}-r_1c_{\theta1})$ por Omega da potencia al fluido. Dividir por caudal produce Euler. Si se define trabajo positivo extraído en turbina se usa el negativo del trabajo al fluido.
\item La salida de turbina debe superar la presión requerida para entrar en cámara con pérdidas posteriores. Esa restricción limita expansión disponible y se acopla a presión de bomba y potencia de eje.
\item En generador, el gas de turbina sigue ruta auxiliar y su masa cuenta en consumo total. En expansor, el calor de refrigeración alimenta entalpía para la turbina. Se deben dibujar bombas y cierre de eje además de esas rutas distintivas.
\end{enumerate}
\section*{Capítulo 20}
\begin{enumerate}
\item De $\dot m_i=C\sqrt{p_s-p_c}$ se obtiene sensibilidad negativa $-\dot m_i/[2(p_s-p_c)]$. El retardo multiplica una respuesta armónica por una fase; el signo estático no decide el efecto dinámico.
\item El promedio es $\widehat p\widehat q\cos\phi/2$: positivo para cero, nulo para $\pi/2$ y negativo para pi. La contribución neta al modo necesita además integración espacial y pérdidas.
\item Reducir paso o tolerancia y recuperar el mismo resultado comprueba discretización. Comparar con observaciones independientes evalúa si el modelo representa el fenómeno; un cálculo convergido puede ser físicamente insuficiente.
\item $I_{sp,tot}=(m_1I_{sp,1}+m_2I_{sp,2})/(m_1+m_2)$. Promediar por tiempo sólo sería válido con caudales iguales y constantes, o integrando correctamente su ponderación.
\end{enumerate}
\section*{Capítulo 21}
\begin{enumerate}
\item De $P_j=\eta P_{el}=\dot mv_e^2/2$ sale el caudal; multiplicarlo por $v_e$ da $F=2\eta P_{el}/v_e$. El impulso específico de ese haz es $v_e/g_0$. Las unidades son kg/s, N y s.
\item $F=2(0.60)(1500)/20000=0.0900\,\mathrm N$ y $\dot m=4.50\times10^{-6}\,\mathrm{kg/s}$. El flujo cinético recupera 900 W.
\item Aplicar Cauchy--Schwarz a $\int v_xd\dot m$ produce $F_x^2\leq\dot m\int v_x^2d\dot m\leq2\dot mP_j$. Igualdad requiere velocidad uniforme y puramente axial.
\item $k_BT_e=3(1.602176634\times10^{-19})$ J y $T_e=E/k_B\simeq34813.6$ K. No se interpreta 3 eV como 3 K ni se multiplica esa energía otra vez por $k_B$.
\end{enumerate}
\section*{Capítulo 22}
\begin{enumerate}
\item $\dot mc_p(T_{out}-T_{in})=\eta_hP_{el}$ da el despeje indicado. Si eta ya incluye calor a pared, no se resta ese calor otra vez del término absorbido.
\item El aumento $\Delta T$ se duplica al reducir caudal a la mitad con los otros parámetros fijos. La temperatura absoluta no se duplica necesariamente porque conserva $T_{in}$.
\item Una descarga puede calentar gas cuya expansión gasdinámica domina aceleración. La clasificación depende del mecanismo de transferencia a movimiento, no de la mera presencia de corriente o ionización.
\item Entrada: $\dot mh_{in}+P_{el}$. Salida: $\dot m(h_{out}+v_e^2/2)$ más energía rechazada y consumos definidos. La entalpía inicial puede contribuir a escape y debe conservar la misma referencia.
\end{enumerate}
\section*{Capítulo 23}
\begin{enumerate}
\item $qV=m_iv_i^2/2$ da rapidez. La corriente cuenta $I_b/q$ iones/s, por lo que $\dot m_i=I_bm_i/q$ y $F=I_b\sqrt{2m_iV/q}$ en el haz axial ideal.
\item Con $F=\lambda\dot m_iv_i$, $\dot m_{tot}=\dot m_i/\eta_u$ y $P_b=\dot m_iv_i^2/2=\eta_eP_{el}$, sustituir en $F^2/(2\dot m_{tot}P_{el})$ da $\lambda^2\eta_u\eta_e$.
\item La corriente de descarga contiene contribuciones que no son iones acelerados del haz. Convertirla íntegra mediante $m_i/q$ sobreestima caudal iónico y altera utilización y empuje inferidos. Se deben separar portadores y cargas.
\item E realiza trabajo $q\mathbf E\cdot\mathbf v$. B magnetiza principalmente electrones y organiza deriva transversal; $q(\mathbf v\times\mathbf B)\cdot\mathbf v=0$. En Hall los iones aceleran sobre todo por el potencial eléctrico.
\end{enumerate}
\section*{Capítulo 24}
\begin{enumerate}
\item $\overline F=fJ_b=40(12\times10^{-6})=4.80\times10^{-4}\,\mathrm N$. Hz es s$^{-1}$ y cancela el segundo del impulso.
\item $I_{sp}=\sum J_i/(g_0\sum m_i)=\sum m_iI_{sp,i}/\sum m_i$. Las masas ponderan cada pulso; la frecuencia sólo cambia el reparto temporal si sus propiedades permanecen.
\item $\Delta E=C(V_i^2-V_f^2)/2$. La potencia media de pulsos idénticos es $f\Delta E$ antes de añadir pérdidas y auxiliares; la instantánea requiere evolución temporal de tensión y corriente.
\item La cuasineutralidad reduce carga neta, no corriente. Cargas opuestas pueden moverse de forma distinta y dar j. La energía eléctrica entra mediante $\mathbf j\cdot\mathbf E$ y se reparte entre movimiento y otras formas.
\end{enumerate}
\section*{Capítulo 25}
\begin{enumerate}
\item El balance material da $m\,dv+v_e\,dm=0$ a primer orden. Con $dm<0$ y escape hacia atrás, $dv>0$. Incluir presión mediante c conserva el mismo signo.
\item $m_p/m_0=1-e^{-0.2}=0.181269$. La aproximación 0.2 excede el valor exacto en 0.018731 absolutos, aproximadamente 10.33\,\% respecto a la fracción exacta.
\item $m_0=m_fe^{\Delta v/c}$ y $m_p=m_0-m_f=m_f(e^{\Delta v/c}-1)$. La masa final incluye carga y estructura unidas; no es sólo carga útil salvo que así se haya definido físicamente.
\item La capacidad ideal no incorpora gravedad, resistencia, dirección, posición final ni duración. Una órbita circular necesita rapidez tangencial a un radio dado; no basta con igualar un presupuesto escalar a $\sqrt{\mu/r}$.
\end{enumerate}
\section*{Capítulo 26}
\begin{enumerate}
\item Las masas son las cuatro sumas del capítulo. En particular $m_{f1}=s_1+m_{02}$. La separación ocurre después de consumir $p_1$ y antes de iniciar el consumo de segunda; no entra como propulsante en ningún logaritmo.
\item $\lambda>0$ requiere $R_m>1$ y $1-\epsilon R_m>0$, de modo que $1<R_m<1/\epsilon$. Numerador nulo significa carga superior tendente a cero para la fracción estructural fijada.
\item Para c y epsilon iguales, la condición de derivadas marginales iguales exige iguales $e^{x_i/c}$ y por tanto iguales $x_i$, en un óptimo interior. La simetría y convexidad del coste apoyan esa solución del modelo.
\item Durante consumo simultáneo, la masa total responde a suma de caudales y el empuje a suma vectorial. Las masas se comparten entre sistemas: no hay dos historias en serie independientes a las que aplicar los logaritmos por separado.
\end{enumerate}
\section*{Capítulo 27}
\begin{enumerate}
\item Integrar $\dot v=-c\dot m/m-g$ con c y g fijos da $c\ln(m_0/m_f)-gt_b$. Si se conserva relación de masas y se duplica duración, la contribución gravitatoria de este modelo se duplica.
\item La forma general contiene $\int g(r)\sin\gamma_fdt$. Gravedad y ángulo varían; una trayectoria horizontal o descendente tiene otra proyección. El término no es disipación de energía mecánica.
\item Presión estática es un estado termodinámico; $q=\rho V_{rel}^2/2$ es una escala cinética. Durante ascenso velocidad puede crecer y densidad caer, por lo que su producto puede tener un máximo interior.
\item Comprobar energía específica y momento angular constante en un campo central sin otras fuerzas. Además se estudia convergencia con tolerancia y cierre de masa; esos controles no equivalen a validación con telemetría.
\end{enumerate}
\section*{Capítulo 28}
\begin{enumerate}
\item Navegación estima el estado; guiado fija la maniobra o referencia compatible con objetivos; control aplica fuerzas y momentos para seguirla. Sus errores y límites se acoplan, pero sus funciones no se intercambian.
\item $|\mathbf v_2-\mathbf v_1|^2=2v^2(1-\cos\Delta i)=4v^2\sin^2(\Delta i/2)$. Para el ángulo usual entre cero y pi se toma raíz positiva: $2v\sin(\Delta i/2)$.
\item $\mathbf M=\mathbf r_T\times\mathbf F_T$ respecto al centro de masas. Si éste se desplaza, cambia el brazo aunque la posición física del actuador permanezca; también puede cambiar inercia.
\item La gravedad actúa durante la maniobra y el vehículo cambia de posición y dirección mientras acelera. Se integra $\dot{\mathbf r}=\mathbf v$, el balance vectorial de fuerzas y $\dot m=-\dot m_p$, con potencia y orientación disponibles.
\end{enumerate}
